\chapter{Własności rozwiązań SDE, procesy Markowa i aproksymacje}
\label{ch:sde-properties}
\paragraph{Odsyłacze międzytomowe.} Estymację parametrów procesów zależnych i dobór modeli omawia tom IV, rozdziały~\ref{IV-ch:time-series} i~\ref{IV-ch:model-selection}.


\section{Od równania do ewolucji prawdopodobieństwa}
W poprzednim rozdziale skonstruowaliśmy proces $X$ na tej
samej przestrzeni co zadany proces Wienera. Chcemy teraz
zrozumieć jego \emph{rozkład} i sposób, w jaki zależy od
warunku początkowego. W tym celu potrzebujemy dwóch
perspektyw. Pierwsza śledzi trajektorie z użyciem wspólnego
szumu; druga opisuje ewolucję wartości oczekiwanych
funkcji testowych. Ich punktem wspólnym jest wzór Itô.

Wszystkie twierdzenia tej sekcji dotyczą rozwiązania
\[
X_t^x=x+\int_0^t b(X_s^x)\,\dd s
+\int_0^t\sigma(X_s^x)\,\dd W_s,\qquad t\ge0,
\tag{6.1}\label{eq:timehom-sde}
\]
gdzie $b:\R^d\to\R^d$, $\sigma:\R^d\to\R^{d\times m}$
są globalnie Lipschitzowskie, a $W$ jest $m$-wymiarowym
procesem Wienera względem filtracji spełniającej warunki
zwykłe. Na każdym skończonym przedziale proces jest
skonstruowany przez twierdzenie \ref{th:sde-global};
jego rozszerzenia na rosnące przedziały są zgodne
dzięki jednoznaczności ścieżkowej.

\section{Własność Markowa wynikająca z jednoznaczności silnej}
\begin{definicja}[Półgrupa przejścia]
Dla $t\ge0$, $x\in\R^d$ i ograniczonej funkcji
borelowskiej $f:\R^d\to\R$ oznaczamy
\[
P_tf(x)=\E[f(X_t^x)],\qquad
P_t(x,A)=\Prob(X_t^x\in A),\quad A\in\calB(\R^d).
\tag{6.2}\label{eq:sde-semigroup}
\]
Wartości te nie zależą od konkretnej przestrzeni
probabilistycznej, na której skonstruowano rozwiązanie:
rozwiązania otrzymane metodą Picarda z identycznych
danych mają te same rozkłady.
\end{definicja}

\begin{lemat}[Zależność rozkładu od warunku początkowego]
\label{lem:sde-kernel}
Dla każdego $t\ge0$ funkcja $x\mapsto P_tf(x)$
jest ciągła, jeżeli $f:\R^d\to\R$ jest ograniczona
i ciągła. Ponadto dla każdej ograniczonej funkcji
borelowskiej $f$ odwzorowanie $x\mapsto P_tf(x)$
jest borelowskie.
\end{lemat}
\begin{proof}
Niech $x_n\to x$ i sprzęgnijmy rozwiązania
$X^{x_n},X^x$ przez ten sam proces $W$.
Z twierdzenia \ref{th:sde-stability} otrzymujemy
$\E\sup_{s\le t}\norm{X_s^{x_n}-X_s^x}^2\to0$.
Stąd $X_t^{x_n}\to X_t^x$ według prawdopodobieństwa.
Jeśli $f$ jest ograniczona i ciągła, to
$f(X_t^{x_n})\to f(X_t^x)$ według prawdopodobieństwa,
a z ograniczoności wynika zbieżność wartości oczekiwanych.
Dla ogólnej funkcji borelowskiej rozpatrujemy klasę
ograniczonych funkcji $f$, dla których $P_tf$ jest
borelowska. Zawiera ona ograniczone funkcje ciągłe;
pozostaje też zamknięta na ograniczone granice punktowe,
ponieważ zbieżność zdominowana działa dla każdej
ustalonej wartości $x$. Ograniczone funkcje ciągłe
wyznaczają borelowskie $\sigma$-ciało na $\R^d$.
Twierdzenie o klasie monotonicznej, zastosowane do
funkcji, daje zatem wszystkie ograniczone funkcje
borelowskie. Jest to zarazem uzasadnienie, że
$P_t(x,A)$ jest jądrem probabilistycznym, a nie
tylko formalnym symbolem.
\end{proof}

\begin{twierdzenie}[Własność Markowa i równania Chapmana--Kołmogorowa]
\label{th:sde-markov}
Niech $X$ będzie rozwiązaniem \eqref{eq:timehom-sde}
z $\calF_0$-mierzalnym $X_0\in L^2$, a $W$ procesem
Wienera względem tej filtracji. Wówczas dla dowolnych
$s,t\ge0$ i ograniczonej funkcji borelowskiej $f$
\[
\E[f(X_{s+t})\mid\calF_s]=P_tf(X_s)
\quad\text{prawie na pewno}.
\tag{6.3}\label{eq:sde-markov}
\]
W szczególności $P_0f=f$ i $P_sP_tf=P_{s+t}f$.
Własność ta dotyczy \emph{całej przeszłości} zawartej
w $\calF_s$, a nie tylko pojedynczej zmiennej $X_s$.
\end{twierdzenie}
\begin{proof}
Ustalmy $s$ i rozważmy szum przesunięty
$\widehat W_u=W_{s+u}-W_s$. Jest on procesem Wienera
niezależnym od $\calF_s$. Na przedziale $[s,s+t]$
proces $Z_u=X_{s+u}$ spełnia równanie
\[
Z_u=X_s+\int_0^u b(Z_v)\,\dd v
+\int_0^u\sigma(Z_v)\,\dd\widehat W_v.
\]
Porównajmy go z iteracjami Picarda uruchomionymi
w chwili $s$: $Z^{(0)}_u=X_s$ i
$Z^{(n+1)}_u=X_s+\int_0^u b(Z_v^{(n)})\,\dd v
+\int_0^u\sigma(Z_v^{(n)})\,\dd\widehat W_v$.
Dla deterministycznego $x$ analogiczne iteracje
zaczynają się w $x$ i korzystają wyłącznie z
$\widehat W$. Dla dowolnej ograniczonej funkcji
borelowskiej $f$ wartość warunkowa
$\E[f(Z_t^{(n)})\mid\calF_s]$ zależy wyłącznie od
$X_s$: w kroku indukcyjnym korzystamy z niezależności
całego procesu $\widehat W$ od $\calF_s$ oraz z
mierzalności iteracji względem $X_s$ i $\widehat W$.
Dokładniej, dla skończenie wielu prostych przybliżeń
całek Itô jest to zwykła własność warunkowania
funkcji dwóch niezależnych zmiennych losowych;
przejście do granicy w $L^2$ uzasadnia tę samą
własność dla iteracji Picarda.

Z twierdzenia \ref{th:sde-global} iteracje zbiegają do
$Z$ w $L^2(\Omega;C([0,t];\R^d))$.
Dla ograniczonego ciągłego $f$ wynika stąd
zbieżność $f(Z_t^{(n)})\to f(Z_t)$ w $L^1$.
Odpowiednie prawa dla deterministycznego $x$
zbiegać będą do $P_tf(x)$; zatem
$\E[f(Z_t)\mid\calF_s]=P_tf(X_s)$.
Przejście do ograniczonych funkcji borelowskich
uzyskujemy przez klasę monotoniczną, używając
lematu \ref{lem:sde-kernel} dla mierzalności
prawej strony. Na końcu korzystamy z tower property:
$P_{s+t}f(x)=\E[P_tf(X_s^x)]=P_s(P_tf)(x)$.
\end{proof}

\begin{twierdzenie}[Mocna własność Markowa]
\label{th:sde-strongmarkov}
Przy założeniach twierdzenia \ref{th:sde-markov}
niech $\tau$ będzie prawie na pewno skończonym
czasem zatrzymania. Dla każdego $t\ge0$ i dowolnej
ograniczonej funkcji borelowskiej $f$ zachodzi
\[
\E[f(X_{\tau+t})\mid\calF_\tau]=P_tf(X_\tau)
\quad\text{prawie na pewno}.
\tag{6.4}\label{eq:sde-strongmarkov}
\]
\end{twierdzenie}
\begin{proof}
Najpierw niech $\tau$ przyjmuje skończenie wiele
wartości $s_1,\ldots,s_N$. Mnożąc
\eqref{eq:sde-markov} przez wskaźnik
$\ind{\{\tau=s_k\}}\in\calF_{s_k}$ i sumując,
otrzymujemy warunkową równość dla $\tau$.

Następnie niech $\tau\le R$ prawie na pewno.
Przybliżamy ją z góry czasami
$\tau_n=2^{-n}\lceil2^n\tau\rceil$,
które są czasami zatrzymania dla filtracji
prawostronnie ciągłej, przyjmują skończenie wiele
wartości i spełniają $\tau_n\downarrow\tau$.
Dla $A\in\calF_\tau$ mamy $A\in\calF_{\tau_n}$,
zatem
\[
\E[\ind A f(X_{\tau_n+t})]
=\E[\ind A P_tf(X_{\tau_n})].
\]
Jeżeli $f$ jest ograniczona i ciągła, to lemat
\ref{lem:sde-kernel}, ciągłość trajektorii $X$
i zbieżność zdominowana pozwalają przejść do granicy,
dając identyczną równość z $\tau$.
Dla ograniczonej funkcji borelowskiej $f$ stosujemy
ponownie twierdzenie o klasie monotonicznej.
Ograniczenie $\tau\le R$ usuwamy przez zastosowanie
wyniku do $\tau\wedge R$ i przejście $R\to\infty$;
na zdarzeniu $\{\tau\le R\}$ odpowiednie tożsamości
pokrywają się, zaś $\Prob(\tau>R)\to0$.
\end{proof}

\begin{przyklad}[Co proces Markowa naprawdę „pamięta”?]
Dla procesu Ornsteina--Uhlenbecka z zadania
\ref{zad:ou}, przy ustalonych $s,t\ge0$,
\[
\E[X_{s+t}\mid\calF_s]
=m+(X_s-m)e^{-\theta t},
\quad
\Var(X_{s+t}\mid\calF_s)
=\frac{\eta^2}{2\theta}(1-e^{-2\theta t}).
\]
To wynika z rozbicia całki stochastycznej w chwili
$s$ i niezależności przyszłych przyrostów. Znajomość
całej trajektorii do chwili $s$ nie zmienia tych
wzorów po ustaleniu $X_s$. Nie jest to jednak
niezależność kolejnych wartości procesu.
\end{przyklad}

\section{Generator dyfuzji i wzór Dynkina}
\begin{definicja}[Operator generatora]
Dla funkcji $f\in C_c^2(\R^d)$ określamy
\[
Lf(x)=\sum_{i=1}^d b_i(x)\partial_i f(x)
+\frac12\sum_{i,j=1}^d a_{ij}(x)\partial_{ij}f(x),
\qquad a(x)=\sigma(x)\sigma(x)^{\mathsf T}.
\tag{6.5}\label{eq:sde-generator}
\]
Operator $L$ nazywamy \emph{generatorem różniczkowym}
dyfuzji. Nie zakładamy w tej definicji, że funkcja
$f$ leży już w dziedzinie generatora silnej półgrupy
w przestrzeni funkcji ograniczonych.
\end{definicja}

\begin{twierdzenie}[Wzór Dynkina w skończonym horyzoncie]
\label{th:sde-dynkin}
Niech $X^x$ będzie rozwiązaniem \eqref{eq:timehom-sde}
z deterministycznym $x\in\R^d$ i niech
$f\in C_c^2(\R^d)$. Wówczas proces
\[
M_t^f=f(X_t^x)-f(x)-\int_0^tLf(X_s^x)\,\dd s
\tag{6.6}\label{eq:dynkin-mart}
\]
jest kwadratowo całkowalnym martyngałem na każdym
skończonym przedziale. Dla dowolnego ograniczonego
czasu zatrzymania $\tau$ otrzymujemy
\[
\E f(X_\tau^x)=f(x)+\E\int_0^\tau Lf(X_s^x)\,\dd s.
\tag{6.7}\label{eq:dynkin-optional}
\]
\end{twierdzenie}
\begin{proof}
Wzór Itô \ref{th:ito-multi} zastosowany do $f(X_t)$
daje
$\dd f(X_t)=Lf(X_t)\,\dd t+
\nabla f(X_t)\sigma(X_t)\,\dd W_t$.
Ponieważ $f$ ma zwarty nośnik, $\nabla f$ i
$D^2 f$ są ograniczone i zerują się poza pewnym
zbiorem zwartym. Współczynniki $b,\sigma$ są
ciągłe i lokalnie ograniczone, więc $Lf$ oraz
$\nabla f\,\sigma$ są \emph{globalnie ograniczone}.
W szczególności
$\E\int_0^T\norm{\nabla f(X_s)\sigma(X_s)}^2\,\dd s
\le T\sup_x\norm{\nabla f(x)\sigma(x)}^2<\infty$.
Stąd całka stochastyczna jest kwadratowo całkowalnym
martyngałem, co daje \eqref{eq:dynkin-mart}.
Dla ograniczonego $\tau$ stosujemy twierdzenie
o próbkowaniu opcjonalnym, a do skończenie zmiennej
w czasie całki względem $s$ — Fubiniego.
\end{proof}

\begin{wniosek}[Generator jako pochodna średniej]
\label{cor:sde-gen-limit}
Przy założeniach twierdzenia \ref{th:sde-dynkin},
dla każdej funkcji $f\in C_c^2(\R^d)$ i każdego
$x\in\R^d$ zachodzi
\[
\lim_{t\downarrow0}\frac{P_tf(x)-f(x)}{t}=Lf(x).
\tag{6.8}\label{eq:genlimit}
\]
\end{wniosek}
\begin{proof}
Z \eqref{eq:dynkin-mart} po wzięciu oczekiwanych mamy
$(P_tf(x)-f(x))/t=(1/t)\int_0^t\E Lf(X_s^x)\,\dd s$.
Funkcja $Lf$ jest ograniczona i ciągła, a
$X_s^x\to x$ prawie na pewno, gdy $s\downarrow0$.
Z twierdzenia o zbieżności zdominowanej wynika
$\E Lf(X_s^x)\to Lf(x)$.
Średnia Cesàro takiej funkcji również dąży do
$Lf(x)$, co dowodzi tezy.
\end{proof}

\begin{uwaga}[Co mierzy generator?]
Dla bardzo małego $h>0$ mamy w przybliżeniu
$\E[f(X_{t+h})\mid X_t=x]-f(x)\approx hLf(x)$.
Pierwsze pochodne w \eqref{eq:sde-generator}
opisują dryf, a pochodne drugiego rzędu są
konsekwencją wariacji kwadratowej szumu.
To związek z równaniami parabolicznymi, który
będzie rozwijany w rozdziałach o półgrupach i
wzorze Feynmana--Kaca.
\end{uwaga}

\begin{zadanie}[Wyjście dyfuzji z przedziału: generator i warunek brzegowy]
\label{zad:exit-drift}
Niech $a>0$, $\mu\in\R$, $\sigma_0>0$ i niech
$X_t=x+\mu t+\sigma_0W_t$, gdzie $x\in(0,a)$.
Dla $\tau=\inf\{t\ge0:X_t\notin(0,a)\}$ obliczyć
$\Prob_x(X_\tau=a)$.

\emph{Rozwiązanie.} W pierwszym kroku dowodzimy, że
$\tau<\infty$ prawie na pewno. Ustalmy $\delta>0$.
Niezależnie od położenia procesu w $(0,a)$ na początku
przedziału długości $\delta$, zdarzenie
$\{\mu\delta+\sigma_0(W_{s+\delta}-W_s)>a\}$
wymusza opuszczenie $(0,a)$ do chwili $s+\delta$.
Ma ono prawdopodobieństwo $q>0$, niezależne od
$\calF_s$. Iterowane warunkowanie daje
$\Prob_x(\tau>n\delta)\le(1-q)^n$.
W szczególności $\E_x\tau<\infty$ i $X_\tau\in\{0,a\}$.

Generator ma postać $L=\mu\partial_x+
(\sigma_0^2/2)\partial_{xx}$. Poszukujemy funkcji
$u\in C^2([0,a])$ spełniającej
$Lu=0$ na $(0,a)$, $u(0)=0$, $u(a)=1$.
Jeśli $\mu\ne0$, równanie
$u''+(2\mu/\sigma_0^2)u'=0$ po jednokrotnym
scałkowaniu daje
$u'(z)=C\exp(-2\mu z/\sigma_0^2)$;
po uwzględnieniu wartości brzegowych otrzymujemy
\[
u(x)=\frac{1-\exp(-2\mu x/\sigma_0^2)}
{1-\exp(-2\mu a/\sigma_0^2)}.
\]
Dla $\mu=0$ równanie $u''=0$ daje $u(x)=x/a$.

Aby wykazać, że to rzeczywiście szukane prawdopodobieństwo,
stosujemy wzór Itô do $u(X_{t\wedge\tau})$.
Ponieważ $Lu=0$ przed wyjściem, jest to lokalny
martyngał. Jest on jednak ograniczony przez $1$,
więc jest prawdziwym martyngałem i
$\E_xu(X_{t\wedge\tau})=u(x)$.
Dzięki prawie pewnej skończoności $\tau$ i ograniczoności
$u$ możemy zastosować zbieżność zdominowaną, uzyskując
$u(x)=\E_xu(X_\tau)=\Prob_x(X_\tau=a)$.
Nie wyciągamy zatem wyniku z samego formalnego rozwiązania
równania różniczkowego: wskazaliśmy, dlaczego jego
rozwiązanie przedstawia prawdopodobieństwo zdarzenia.
\end{zadanie}

\section{Równanie w przód: rozkład prawdopodobieństwa jako rozwiązanie słabe}
\begin{twierdzenie}[Słabe równanie Fokkera--Plancka]
\label{th:sde-fp}
Niech $X$ będzie rozwiązaniem \eqref{eq:timehom-sde}
z $X_0\in L^2(\Omega;\R^d)$, a
$\mu_t=\Law(X_t)$ dla $t\ge0$. Wówczas dla każdej
$f\in C_c^2(\R^d)$
\[
\int_{\R^d}f(x)\,\dd\mu_t(x)
=\int_{\R^d}f(x)\,\dd\mu_0(x)
+\int_0^t\int_{\R^d}Lf(x)\,\dd\mu_s(x)\,\dd s.
\tag{6.9}\label{eq:fp-weak}
\]
Jeśli ponadto dla $t>0$ rozkład $\mu_t$ ma gęstość
$\rho(t,\cdot)$, funkcja $\rho$ jest klasy $C^1$
w zmiennej $t$, a funkcje $b_i\rho$ i $a_{ij}\rho$
mają ciągłe pochodne przestrzenne odpowiednio rzędu
pierwszego i drugiego, przy lokalnej całkowalności
umożliwiającej różniczkowanie całek funkcji testowych,
to $\rho$ spełnia klasycznie równanie
\[
\partial_t\rho
=-\sum_i\partial_i(b_i\rho)
+\tfrac12\sum_{i,j}\partial_{ij}(a_{ij}\rho),
\qquad a=\sigma\sigma^{\mathsf T}.
\tag{6.10}\label{eq:fp-formal}
\]
Bez tych dodatkowych założeń \eqref{eq:fp-weak}
jest właściwym, matematycznie ścisłym sformułowaniem.
\end{twierdzenie}
\begin{proof}
Ze wzoru Itô i oszacowania w dowodzie twierdzenia
\ref{th:sde-dynkin} wynika
$\E f(X_t)-\E f(X_0)=\E\int_0^t Lf(X_s)\,\dd s$.
Ponieważ $Lf$ jest ograniczone, twierdzenie Fubiniego
daje \eqref{eq:fp-weak} po zapisaniu wartości
oczekiwanych jako całek względem rozkładów $\mu_s$.
W przypadku istnienia regularnej gęstości podstawiamy
$\dd\mu_s(x)=\rho(s,x)\,\dd\lambda_d(x)$.
Dla każdej funkcji testowej o zwartym nośniku możemy
różniczkować lewą stronę względem czasu (na dowolnym
zwartym podprzedziale dodatnich chwil). Po scałkowaniu
przez części w przestrzeni dostajemy
$\int f\,\partial_t\rho\,\dd\lambda_d
=\int f[-\partial_i(b_i\rho)+\frac12\partial_{ij}(a_{ij}\rho)]
\,\dd\lambda_d$ (z sumowaniem po indeksach).
Wszystkie te funkcje są ciągłe na rozpatrywanym obszarze.
Lemat podstawowy rachunku wariacyjnego, stosowany do
$C_c^\infty(\R^d)$, utożsamia ich wartości punktowo.
To uzasadnia postać klasyczną \eqref{eq:fp-formal},
a nie tylko jej zapis formalny.
\end{proof}

\begin{zadanie}[Ruch Browna i równanie ciepła]
Niech $X_t=x+W_t$ w $\R^d$. Wykazać, że jego gęstość
spełnia równanie ciepła z generatorem $\tfrac12\Delta$.

\emph{Rozwiązanie.} Tutaj $b=0$, $\sigma=I_d$,
a więc $Lf=\tfrac12\Delta f$. Dla $t>0$ gęstość
\[
p_t(x,y)=\frac1{(2\pi t)^{d/2}}
\exp\!\left(-\frac{\norm{y-x}^2}{2t}\right)
\]
spełnia $\partial_t p_t=\frac12\Delta_y p_t$:
bezpośrednie różniczkowanie daje
$\partial_t p_t=p_t(-d/(2t)+\norm{y-x}^2/(2t^2))$
oraz
$\Delta_y p_t=p_t(-d/t+\norm{y-x}^2/t^2)$.
Tożsamość jest zatem sprawdzona wprost.
Co ważne, w chwili $t=0$ prawo procesu jest miarą
Diraca $\delta_x$, nie zwykłą funkcją-gęstością.
Stąd naturalne jest słabe sformułowanie
\eqref{eq:fp-weak}.
\end{zadanie}

\begin{zadanie}[Rozkład stacjonarny procesu Ornsteina--Uhlenbecka]
Niech $X$ spełnia równanie z zadania \ref{zad:ou}
ze stałymi $\theta,\eta>0$ i $m\in\R$.
Znaleźć rozkład niezmienniczy i sprawdzić jego
niezmienniczość.

\emph{Rozwiązanie.} Z \eqref{eq:ou-solution}
wynika, że dla deterministycznego $X_0=x$ rozkład
$X_t$ jest normalny o średniej
$m+(x-m)e^{-\theta t}$ i wariancji
$\eta^2(1-e^{-2\theta t})/(2\theta)$.
Kandydatem na rozkład niezmienniczy jest
$\pi=\Normal(m,\eta^2/(2\theta))$.
Jeżeli $X_0\sim\pi$ i $X_0$ jest niezależny od
przyszłych przyrostów $W$ (co wynika z założeń
filtracji), to \eqref{eq:ou-solution} jest sumą
niezależnych zmiennych normalnych, a jej średnia
wynosi $m$ i wariancja
\[
\frac{\eta^2}{2\theta}e^{-2\theta t}
+\frac{\eta^2}{2\theta}(1-e^{-2\theta t})
=\frac{\eta^2}{2\theta}.
\]
Stąd $\Law(X_t)=\pi$ dla każdego $t\ge0$.
To sprawdza niezmienniczość. Zbieżność rozkładów
przy deterministycznym punkcie startowym, otrzymana
w poprzednim rozdziale, wyjaśnia fizyczną intuicję
powrotu do równowagi.
\end{zadanie}

\section{Porównywanie rozwiązań w jednym wymiarze}
\begin{twierdzenie}[Zasada porównania]
\label{th:sde-comparison}
Niech $b_1,b_2,\sigma:[0,T]\times\R\to\R$
będą borelowskie, globalnie Lipschitzowskie w
zmiennej przestrzennej jednostajnie względem czasu
i ograniczone w zerze. Zakładamy, że
$b_1(t,x)\le b_2(t,x)$ dla wszystkich $t,x$.
Niech $X,Y$ będą silnymi rozwiązaniami napędzanymi
tym samym procesem Wienera, odpowiednio,
\[
\dd X_t=b_1(t,X_t)\,\dd t+\sigma(t,X_t)\,\dd W_t,
\quad
\dd Y_t=b_2(t,Y_t)\,\dd t+\sigma(t,Y_t)\,\dd W_t,
\]
z $X_0,Y_0\in L^2(\Omega)$ i $X_0\le Y_0$ prawie
na pewno. Wówczas
$\Prob(X_t\le Y_t\ \text{dla wszystkich }t\in[0,T])=1$.
\end{twierdzenie}
\begin{proof}
Niech $D_t=X_t-Y_t$. Wybierzmy gładką, niemalejącą
funkcję $g:\R\to[0,1]$, dla której $g(z)=0$
przy $z\le0$ oraz $g(z)=1$ przy $z\ge1$.
Istnieje stała $C_g<\infty$ taka, że
$0\le g'(z)\le C_g\ind{\{0<z<1\}}$.
Dla $\varepsilon>0$ określmy
$\psi_\varepsilon(z)=\int_0^{z^+}g(u/\varepsilon)\,\dd u$.
Jest to funkcja klasy $C^2$, wypukła, spełniająca
\[
0\le\psi_\varepsilon'(z)\le\ind{\{z>0\}},\quad
0\le z^+-\psi_\varepsilon(z)\le\varepsilon,\quad
0\le\psi_\varepsilon''(z)\le
\frac{C_g}{\varepsilon}\ind{\{0<z<\varepsilon\}}.
\]
Możemy więc zastosować zwykły wzór Itô bez
odwoływania się do niesformułowanej wersji dla
funkcji z nieciągłą drugą pochodną. Po wzięciu
wartości oczekiwanych
całka stochastyczna znika, ponieważ
$\psi_\varepsilon'$ jest ograniczona, a
$\sigma(t,X_t)-\sigma(t,Y_t)$ jest kwadratowo
całkowalna. Otrzymujemy
\[
\begin{split}
\E\psi_\varepsilon(D_t)
=&\E\psi_\varepsilon(D_0)
+\E\int_0^t\psi_\varepsilon'(D_s)
[b_1(s,X_s)-b_2(s,Y_s)]\,\dd s\\
&+\frac12\E\int_0^t\psi_\varepsilon''(D_s)
[\sigma(s,X_s)-\sigma(s,Y_s)]^2\,\dd s.
\end{split}
\]
Pierwszy składnik jest zerowy, gdyż $D_0\le0$.
Na zbiorze $D_s>0$ mamy
$b_1(s,X_s)-b_2(s,Y_s)
\le b_1(s,X_s)-b_1(s,Y_s)\le L D_s^+$.
Dla wyrazu drugiego rzędu warunek Lipschitza daje
\[
\psi_\varepsilon''(D_s)
[\sigma(s,X_s)-\sigma(s,Y_s)]^2
\le \frac{C_gL^2}{\varepsilon}D_s^2
\ind{\{0<D_s<\varepsilon\}}\le C_gL^2\varepsilon.
\]
Wobec tego
$\E\psi_\varepsilon(D_t)
\le L\int_0^t\E D_s^+\,\dd s+
\tfrac12 C_gL^2\varepsilon t$.
Ponieważ $0\le D_t^+-\psi_\varepsilon(D_t)\le\varepsilon$,
po przejściu $\varepsilon\downarrow0$ mamy
$\E D_t^+\le L\int_0^t\E D_s^+\,\dd s$.
Lemat Gronwalla daje $\E D_t^+=0$ dla każdego $t$.
Na przeliczalnym zbiorze gęstym i dzięki ciągłości
trajektorii otrzymujemy żądaną nierówność
jednocześnie dla wszystkich chwil.
\end{proof}

\begin{uwaga}
Kluczowym założeniem jest \emph{wspólny} współczynnik
dyfuzji $\sigma$, a nie tylko porządek dryfów.
Jeśli różne szumy działają na oba procesy, nawet
początkowy porządek nie musi się utrzymać.
\end{uwaga}

\section{Metoda Eulera--Maruyamy: pierwszy schemat numeryczny}
\begin{definicja}[Schemat Eulera--Maruyamy]
\label{def:em}
Niech $T>0$, $h=T/N$ i $t_k=kh$, $k=0,\ldots,N$.
Dla funkcji $b,\sigma$ z rozdziału
\ref{ch:sde-existence} przyjmujemy
$\bar X_0=X_0$ oraz
\[
\bar X_{k+1}=\bar X_k+b(t_k,\bar X_k)h
+\sigma(t_k,\bar X_k)(W_{t_{k+1}}-W_{t_k}).
\tag{6.11}\label{eq:em}
\]
Ciągły interpolant \emph{stochastyczny} definiujemy
przez $X_t^h=X_0+\int_0^t b(\eta_h(s),X^h_{\eta_h(s)})\,\dd s
+\int_0^t\sigma(\eta_h(s),X^h_{\eta_h(s)})\,\dd W_s$,
gdzie $\eta_h(s)=t_k$ dla $s\in[t_k,t_{k+1})$.
W punktach siatki $X_{t_k}^h=\bar X_k$.
\end{definicja}

\begin{uwaga}[Dwa rodzaje błędu]
Błąd \emph{silny} kontroluje różnicę $X^h-X$
na tej samej realizacji szumu. Błąd \emph{słaby}
dotyczy różnicy wartości oczekiwanych
$\E f(X_T^h)-\E f(X_T)$ dla funkcji testowych.
Poniżej dowodzimy wyłącznie oszacowania silnego;
błędy słabe wymagają odrębnych założeń i analizy.
\end{uwaga}

\begin{twierdzenie}[Silna zbieżność metody Eulera--Maruyamy]
\label{th:em-rate}
Niech $X$ będzie rozwiązaniem SDE z twierdzenia
\ref{th:sde-global}, a $X^h$ interpolantem
\ref{def:em}, napędzanym tym samym procesem Wienera
i z tym samym $X_0\in L^2$. Zakładamy dodatkowo,
że istnieje $K$ takie, iż dla wszystkich $s,t\in[0,T]$
i $x\in\R^d$
\[
\norm{b(t,x)-b(s,x)}+
\norm{\sigma(t,x)-\sigma(s,x)}_{\mathrm F}
\le K(1+\norm{x})|t-s|^{1/2}.
\tag{6.12}\label{eq:time-holder}
\]
Wówczas dla pewnej stałej $C=C(T,L,K,\E\norm{X_0}^2)$,
niezależnej od $N$,
\[
\E\sup_{t\le T}\norm{X_t-X_t^h}^2\le Ch.
\tag{6.13}\label{eq:em-strong}
\]
Oznacza to silny rząd $1/2$ w normie
$\|\sup_{t\le T}\norm{X_t-X_t^h}\|_{L^2}$.
\end{twierdzenie}
\begin{proof}
\emph{Krok 1: jednostajne ograniczenie momentów schematu.}
Z definicji interpolantu, liniowego wzrostu
współczynników i lematu \ref{lem:sde-sup} otrzymujemy
\[
\E\sup_{r\le t}\norm{X_r^h}^2
\le C_0(1+\E\norm{X_0}^2)
+C_1\int_0^t\E\norm{X^h_{\eta_h(s)}}^2\,\dd s.
\]
Ponieważ $\eta_h(s)\le s$, ostatnia wartość
oczekiwana jest nie większa niż
$\E\sup_{u\le s}\norm{X_u^h}^2$.
Z lematu Gronwalla
\[
\sup_{N\ge1}\E\sup_{t\le T}\norm{X_t^h}^2<\infty.
\tag{6.14}\label{eq:em-moment}
\]

\emph{Krok 2: przyrost wewnątrz pojedynczego kroku.}
Dla $s\in[t_k,t_{k+1})$, warunkując względem
$\calF_{t_k}$ i korzystając z izometrii,
otrzymujemy
\[
\begin{split}
\E\norm{X_s^h-X_{t_k}^h}^2
&\le2(s-t_k)^2\E\norm{b(t_k,X_{t_k}^h)}^2\\
&\quad+2(s-t_k)\E\norm{\sigma(t_k,X_{t_k}^h)}_{\mathrm F}^2
\le C_2h,
\end{split}
\]
gdzie ostatni krok wynika z liniowego wzrostu i
\eqref{eq:em-moment}. W szczególności
$\int_0^T\E\norm{X_s^h-X_{\eta_h(s)}^h}^2\,\dd s
\le C_2Th$.

\emph{Krok 3: porównanie równań całkowych.}
Oznaczmy $e_t=X_t-X_t^h$. Odejmując równania,
a następnie stosując lemat \ref{lem:sde-sup},
rozkładamy różnice współczynników na trzy składniki:
\[
\begin{split}
\norm{b(s,X_s)-b(\eta_h(s),X^h_{\eta_h(s)})}
&\le L\norm{X_s-X_s^h}
+L\norm{X_s^h-X^h_{\eta_h(s)}}\\
&\quad+K(1+\norm{X^h_{\eta_h(s)}})\sqrt h,
\end{split}
\]
i analogiczną nierówność dla macierzy $\sigma$.
Po podniesieniu do kwadratu, scałkowaniu i
wykorzystaniu kroków 1--2 otrzymujemy
\[
\E\sup_{r\le t}\norm{e_r}^2
\le C_3\int_0^t\E\sup_{u\le s}\norm{e_u}^2\,\dd s
+C_4h.
\]
Lemat Gronwalla daje \eqref{eq:em-strong}.
Oszacowanie jest jednostajne na całym przedziale,
a nie tylko w punktach siatki.
\end{proof}

\begin{zadanie}[Schemat Eulera dla procesu Ornsteina--Uhlenbecka]
Dla równania $\dd X_t=\theta(m-X_t)\,\dd t+\eta\,\dd W_t$
z deterministycznym $X_0=x_0$ wyznaczyć rekurencję
Eulera, jej średnią i wariancję po $k$ krokach.

\emph{Rozwiązanie.} Ze wzoru \eqref{eq:em}
\[
\bar X_{k+1}=(1-\theta h)\bar X_k+\theta hm
+\eta\Delta_kW.
\]
Wprowadźmy $a=1-\theta h$. Po odjęciu $m$
i iteracji otrzymujemy
$\bar X_k-m=a^k(x_0-m)+\eta\sum_{j=0}^{k-1}
a^{k-1-j}\Delta_jW$. Przyrosty są niezależne,
mają wariancję $h$, więc
\[
\E\bar X_k=m+a^k(x_0-m),\qquad
\Var(\bar X_k)=\eta^2h\sum_{j=0}^{k-1}a^{2j}.
\]
Jeżeli $|a|<1$, to wariancja graniczna schematu
przy stałym $h$ wynosi $\eta^2h/(1-a^2)$,
co na ogół nie jest równe dokładnej wariancji
stacjonarnej $\eta^2/(2\theta)$.
To pokazuje, dlaczego dyskretyzacja może zmieniać
własności długoczasowe, nawet gdy zbiega na każdym
ustalonym przedziale wraz z $h\downarrow0$.
\end{zadanie}

\begin{center}
\begin{tikzpicture}[>=Latex,x=1cm,y=.7cm]
\draw[->] (0,0)--(10,0) node[right] {$t$};
\draw[->] (0,-.3)--(0,3) node[above] {$X_t$};
\draw[thick,blue!70!black] (0,.6)--(.6,.95)--(1.2,.72)--(1.8,1.48)
--(2.4,1.3)--(3,1.65)--(3.6,1.12)--(4.2,1.85)
--(4.8,1.43)--(5.4,1.7)--(6,2.35)--(6.6,1.9)
--(7.2,2.55)--(7.8,2.06)--(8.4,2.53)--(9,2.18);
\draw[very thick,orange!85!black] (0,.6)--(1.8,1.25)--(3.6,1.05)
--(5.4,1.65)--(7.2,2.32)--(9,2.2);
\foreach \x in {0,1.8,3.6,5.4,7.2,9} \draw[dashed,gray!55] (\x,0)--(\x,2.8);
\node[blue!70!black,anchor=west,scale=.82] at (5.8,2.95)
{trajektoria procesu (schemat)};
\node[orange!75!black,anchor=west,scale=.82] at (5.9,.6)
{przybliżenie na siatce};
\end{tikzpicture}
\par\smallskip\small Ilustracja idei dyskretyzacji; nie przedstawia obliczonej próbki konkretnego SDE.
\end{center}

\section{Kryterium nieeksplozji poza przypadkiem liniowego wzrostu}
\begin{twierdzenie}[Przykład kryterium Lapunowa]
\label{th:lyapunov-example}
Niech $W$ będzie jednowymiarowym procesem Wienera,
a $X$ lokalnym rozwiązaniem
\[
\dd X_t=-X_t^3\,\dd t+\dd W_t,\qquad X_0=x\in\R.
\]
Współczynniki są lokalnie Lipschitzowskie, lecz dryf
nie ma liniowego wzrostu. Mimo tego rozwiązanie jest
nieeksplozywne na każdym skończonym przedziale.
\end{twierdzenie}
\begin{proof}
Z obcięcia współczynników, jak w dowodzie twierdzenia
\ref{th:sde-local}, otrzymujemy zgodne lokalne
rozwiązania do czasów $\tau_R=\inf\{t:\abs{X_t}\ge R\}$.
Stosujemy wzór Itô do $V(x)=x^2$ i procesu
zatrzymanego w chwili $\tau_R$:
\[
\abs{X_{t\wedge\tau_R}}^2
=x^2+\int_0^{t\wedge\tau_R}(1-2X_s^4)\,\dd s
+2\int_0^{t\wedge\tau_R}X_s\,\dd W_s.
\]
Na przedziale przed $\tau_R$ integrand ostatniej
całki jest ograniczony przez $R$, a więc całka
ma średnią zero. Dostajemy
$\E\abs{X_{t\wedge\tau_R}}^2\le x^2+t$.
Na zdarzeniu $\{\tau_R\le t\}$ ciągłość trajektorii
daje $\abs{X_{t\wedge\tau_R}}=R$, zatem
$R^2\Prob(\tau_R\le t)\le x^2+t$.
Po przejściu $R\to\infty$ wnioskujemy, że czas
życia lokalnego rozwiązania jest nieskończony
prawie na pewno.
\end{proof}

\begin{uwaga}
Powyższy argument jest przykładem metody funkcji
Lapunowa. Ujemny składnik $-2X^4$ przeciwdziała
ucieczce rozwiązania do nieskończoności. W ogólnej
wersji poszukuje się nieujemnej funkcji $V$, która
rośnie przy $\norm{x}\to\infty$ i spełnia odpowiednią
nierówność dla generatora $LV$.
\end{uwaga}

\section{Dwa przykłady pokazujące granice twierdzeń}
\begin{przyklad}[Silna i słaba jednoznaczność dla równania deterministycznego]
Równanie $\dd X_t=2\sqrt{|X_t|}\,\dd t$, $X_0=0$,
ma zarówno rozwiązanie $X_t\equiv0$, jak i
$X_t=(t-c)_+^2$ dla każdego $c\ge0$.
Możemy nadać dowolne rozkłady losowemu opóźnieniu $c$,
otrzymując różne prawa na przestrzeni trajektorii.
Nie ma zatem ani jednoznaczności ścieżkowej, ani
jednoznaczności co do prawa. Oba pojęcia są różne,
choć w tym przykładzie zawodzą jednocześnie.
\end{przyklad}
\begin{przyklad}[Dlaczego sama zbieżność rozkładów nie daje zbieżności silnej?]
Niech $W$ i $\widehat W$ będą niezależnymi ruchami Browna
na tej samej przestrzeni. Dla dowolnego $t>0$ mają identyczny
rozkład $\Normal(0,t)$, ale
$\E|W_t-\widehat W_t|^2=2t$.
Nie można więc wnioskować o błędzie silnym metody
numerycznej ze zbieżności jej rozkładów jednowymiarowych.
W twierdzeniu \ref{th:em-rate} proces i przybliżenie
korzystają z \emph{tych samych} przyrostów Wienera.
\end{przyklad}

\section{Podsumowanie i dalsza droga}
W rozdziałach 5--6 zbudowaliśmy rozwiązania SDE i
wyprowadziliśmy ich podstawowe własności bez
zakładania z góry teorii rozwiązań. Globalny warunek
Lipschitza zapewnia istnienie silne i jednoznaczność
ścieżkową, a lokalny warunek wraz z odpowiednią
kontrolą wzrostu — brak eksplozji. Własność Markowa
i generator wiążą proces z równaniami różniczkowymi
cząstkowymi, natomiast zasada porównania i metoda
Eulera--Maruyamy opisują zależność od danych i
możliwość obliczeń numerycznych.

W dalszej części tomu przejdziemy do słabych
rozwiązań SDE, problemów martyngałowych, wzoru
Feynmana--Kaca, zmiany miary Girsanowa oraz
bardziej ogólnego rachunku dla półmartyngałów.

% ============================================================================
% ROZDZIAL 7 — Problem martyngalowy i Feynman--Kac
% ============================================================================
